\item El cerebro humano y la médula espinal están sumergidos en el fluido cerebroespinal. El fluido normalmente es continuo entre las cavidades craneal y espinal y ejerce una presión de $100$ a $200 \text{ mm de H}_2\text{O}$ sobre la presión atmosférica prevaleciente. En el trabajo médico, las presiones usualmente se miden en milímetros de $\text{H}_2\text{O}$ porque los fluidos corporales, incluido el fluido cerebroespinal, por lo general tienen la misma densidad que el agua. La presión del fluido cerebroespinal se puede medir mediante una sonda espinal, como se ilustra en la figura P14.8. Un tubo hueco se inserta en la columna vertebral y se observa la altura a la que se eleva el fluido. Si el fluido se eleva a una altura de $160 \text{ mm}$, su presión manométrica se escribe como $160 \text{ mm H}_2\text{O}$. (a) Exprese esta presión en pascales, en atmósferas y en milímetros de mercurio. (b) Algunas condiciones que bloquean o inhiben el flujo de fluido cerebroespinal se pueden investigar mediante la prueba de Queckenstedt. En este procedimiento, se comprimen las venas en la nuca del paciente para hacer que la presión sanguínea se eleve en el cerebro, lo cual se transmite al fluido cerebroespinal. Explique cómo el nivel de fluido en la sonda espinal puede utilizarse como herramienta de diagnóstico para la condición de la espina del paciente.
